% CHƯƠNG 6 - Rubric: Đánh giá so sánh với các công nghệ/thuật toán khác
\chapter{Đánh giá so sánh}
\label{chuong:sosanh}

\section{So sánh với dự báo không dùng FGL}

Phép so sánh cốt lõi của phương pháp là giữa Baseline (không dùng teacher) và Student (có FGL),
cùng kiến trúc, cùng dữ liệu huấn luyện, khác nhau duy nhất ở việc có hay không có thành phần
chưng cất tri thức trong hàm mất mát. Trên Mackey--Glass, kết quả ở
Chương~\ref{chuong:taihien} cho thấy FGL cải thiện \textbf{nhất quán và rõ rệt}: toàn bộ sáu
seed của kịch bản chuẩn đều cho Student MSE thấp hơn Baseline MSE, mức giảm trung bình
26,42\% (độ lệch chuẩn 3,86 điểm phần trăm); ở thực nghiệm quét tầm dự báo với Page--Hinkley,
13 trên 14 giá trị horizon cũng cho cải thiện, trung bình 12,22\% (loại horizon=2). Trên
CHB-MIT (ba bệnh nhân 1, 3, 5, Bảng~\ref{bang:chbmit-sosanh-benhnhan}), mức tăng trung bình chỉ
$+31,2\%$ --- thấp hơn 44,8\% bài báo công bố --- và không đồng đều như Mackey--Glass: gần như
không cải thiện ở bệnh nhân 1 (Baseline đã mạnh), nhưng cải thiện mạnh ở bệnh nhân 3 và 5
(Baseline yếu). Phân tích thêm ở mục~\ref{sec:nhanxet-phephan}.

\section{So sánh với các hướng tiếp cận liên quan}

FGL giải quyết vấn đề thay đổi theo thời gian của dữ liệu một cách \textit{ngầm định}: không
có bước phát hiện tường minh thời điểm phân phối dữ liệu thay đổi, mà dựa hoàn toàn vào việc
teacher --- huấn luyện trên dữ liệu gần đích hơn --- truyền tri thức cho student qua chưng cất.
Cách tiếp cận này khác với các phương pháp phát hiện trôi khái niệm (concept drift) tường
minh, chẳng hạn kiểm định Page--Hinkley \cite{page1954cusum} mà chính tác giả cũng thử nghiệm
kết hợp thêm trong \texttt{drift\_exp.py}: ở đó, hệ thống chủ động theo dõi độ lệch giữa giá
trị quan sát và trung bình trượt của nó, và khi độ lệch tích lũy vượt ngưỡng thì trigger huấn
luyện lại student ngay trong lúc suy luận. Hai cơ chế này không loại trừ nhau --- kết quả ở
Bảng~\ref{bang:mg-drift} cho thấy có thể kết hợp cả hai (FGL \textit{và} Page--Hinkley) --- và
việc kết hợp này do chính tác giả thực hiện, không phải đóng góp riêng của báo cáo. Điểm cần
lưu ý là thực nghiệm kết hợp này dùng $\alpha=0$ (student học hoàn toàn từ teacher, không dùng
nhãn thật), nên đo một khía cạnh khác với thực nghiệm chuẩn ($\alpha=0,5$,
mục~\ref{sec:mg-sweep-seed}) --- hai con số 26,42\% và 12,22\% không nên gộp chung so sánh như
cùng một cấu hình.

So với các phương pháp chưng cất tri thức cổ điển \cite{hinton2015distillation}, nơi teacher
và student thường giải \textit{cùng} một bài toán (ví dụ cùng phân loại ảnh) và khác nhau chủ
yếu ở dung lượng mô hình (teacher lớn, student nhỏ), đóng góp đặc thù của FGL là teacher và
student giải \textit{hai bài toán khác nhau về thời điểm} (như phân tích ở
mục~\ref{sec:cochetuonglai}): sự khác biệt về hiệu năng không đến từ việc một mô hình lớn hơn
mô hình khác, mà đến từ việc một mô hình được cho nhìn vào dữ liệu gần đích dự báo hơn.

\section{Nhận xét phê phán}
\label{sec:nhanxet-phephan}

Một số điểm trong kết quả tái hiện chưa hoàn toàn ủng hộ các tuyên bố tổng quát của bài báo,
hoặc còn cần thêm bằng chứng để khẳng định chắc chắn:

\begin{itemize}
	\item \textbf{Ngoại lệ ở tầm dự báo ngắn}: trong Bảng~\ref{bang:mg-drift}, horizon=2 là
	trường hợp duy nhất FGL làm \textit{tệ hơn} baseline ($-7,20\%$). Điều này cho thấy lợi ích
	của FGL không phải tuyệt đối ở mọi cấu hình, mà có thể phụ thuộc vào việc bài toán của
	student đã đủ ``khó'' để teacher có giá trị dẫn dắt hay chưa --- với tầm dự báo rất ngắn,
	bài toán của student đã gần với bài toán của teacher, nên lợi ích chưng cất giảm hoặc mất
	tác dụng.
	\item \textbf{Mức cải thiện tái hiện cao hơn mức công bố}: trung bình 26,42\% (sáu seed) so
	với 23,4\% công bố. Chênh lệch không lớn và cùng hướng, nhưng báo cáo không có đủ cơ sở để
	khẳng định nguyên nhân chính xác của sự khác biệt này (khác biệt về cách loại outlier, hay
	về số lượng cấu hình được lấy trung bình trong bài báo), nên chỉ nêu đây là một quan sát
	cần làm rõ thêm, không khẳng định tái hiện vượt trội hơn bài báo.
	\item \textbf{Độ tin cậy thống kê còn hạn chế}: mỗi seed chỉ được chạy một lần, nên chưa
	tách hoàn toàn được phương sai do khởi tạo trọng số ngẫu nhiên khỏi phương sai do seed chia
	dữ liệu. Kết quả nên được hiểu là ``khớp hướng và độ lớn'' với bài báo, chưa phải một
	khẳng định thống kê chặt (ví dụ kiểm định giả thuyết với nhiều lần lặp cho mỗi seed).
	\item \textbf{Mức tăng AUC-ROC trên CHB-MIT thấp hơn công bố và không đồng đều giữa bệnh
	nhân}: bài báo báo cáo mức tăng trung bình 44,8\%, tái hiện trên ba bệnh nhân chỉ được
	$+31,2\%$ (Bảng~\ref{bang:chbmit-sosanh-benhnhan}), nguyên nhân chính là Baseline hoạt động
	rất khác nhau giữa bệnh nhân (AUC từ 0,20 đến 0,92) --- nơi Baseline đã mạnh (bệnh nhân 1)
	thì FGL gần như không còn gì để cải thiện. Một phần khác biệt cũng có thể do mã nguồn
	CHB-MIT không cố định seed ngẫu nhiên và báo cáo mới tái hiện được ba bệnh nhân với vài
	trial mỗi bệnh nhân, trong khi bài báo có thể đã lấy trung bình trên nhiều bệnh nhân và lần
	chạy hơn. Đây là điểm tái hiện \textbf{không ủng hộ rõ} tuyên bố của bài báo trên bộ dữ liệu
	này, cần nêu rõ thay vì bỏ qua.
	\item \textbf{Chỉ số AUC bất thường ở lần chạy local đã được làm rõ thêm}: AUC-ROC = 0,28 của
	teacher (bệnh nhân 1) trên máy cá nhân không lặp lại ở bất kỳ lần chạy nào trên Kaggle
	(mục~\ref{sec:chbmit-ketqua}) --- nhiều khả năng là một lần khởi tạo ngẫu nhiên không may,
	hơn là lỗi hệ thống. Nghi vấn \texttt{shuffle=False} (mục~\ref{item:shuffle-false}) vẫn là
	một quan sát hợp lệ về thiết kế mã nguồn, nhưng không còn đủ cơ sở khẳng định là nguyên
	nhân chính của riêng kết quả 0,28 đó.
\end{itemize}
